\documentclass[11pt]{article}

\usepackage[margin=1in]{geometry}
\usepackage{amsmath,amssymb,amsthm,mathtools}
\usepackage{booktabs}
\usepackage{microtype}
\usepackage{enumitem}
\usepackage{hyperref}

\hypersetup{colorlinks=true,linkcolor=blue,citecolor=blue,urlcolor=blue}

\newtheorem{theorem}{Theorem}[section]
\newtheorem{proposition}[theorem]{Proposition}
\newtheorem{lemma}[theorem]{Lemma}
\newtheorem{corollary}[theorem]{Corollary}
\newtheorem{example}[theorem]{Example}
\theoremstyle{definition}
\newtheorem{definition}[theorem]{Definition}
\theoremstyle{remark}
\newtheorem{remark}[theorem]{Remark}

\newcommand{\Qx}{\mathcal Q_x}
\newcommand{\Dx}{\mathcal D_x}
\newcommand{\ZZ}{\mathbb Z}
\newcommand{\CC}{\mathbb C}
\newcommand{\one}{\mathbf 1}
\newcommand{\ee}{\mathrm e}
\newcommand{\Edir}{\mathcal E_{\mathrm{dir}}}
\newcommand{\norm}[1]{\left\lVert #1\right\rVert}

\title{\textbf{Direct-Sum Row-Energy Envelope and the Cauchy Bottleneck}}
\author{Liang Wang\\
Huazhong University of Science and Technology\\
Wuhan 430074, P.R. China}
\date{August 2026}

\begin{document}
\maketitle

\begin{abstract}
The transition-native twin-prime program contains a complete-period direct-sum
row energy which remains after TPC-215 controls reduced-frequency clustering.
This paper gives a deterministic bound for that quantity on the literal V46
scales.  Writing the emitter as a sum over primes, the condition
\(4Q<H\) makes the integer cutoff injective modulo each divisor for every fixed
prime.  The fixed-prime row norm is therefore exact, and one Cauchy inequality
across the prime shell yields
\[
 L^{-1}\Edir \ll_{\psi} (Q^3/H)(\log U)^3
 = x^{11/32+o(1)}.
\]
No prime-number theorem and no Mobius cancellation are used.  An exact rational
fixture with all prime rows supported on the same two residues shows that the
shell Cauchy factor cannot be replaced by free orthogonality.  The result is a
complete-period structural envelope only: the finite physical-window Gram,
prime-shell reassembly, arithmetic \(L^2\), and the twin-prime endpoint remain
open.
\end{abstract}

\paragraph{Claim level.}
\texttt{PROVED\_STRUCTURAL\_L1 / DIRECT\_SUM\_ROW\_ENERGY\_ENVELOPE}.\\
The exponent \(11/32\) is an unconditional source-locked envelope.  The finite
fixture is exact rational structural QA.  No arithmetic saving is claimed.

\section{Introduction}

The V46 transition scalar uses the literal coefficient
\begin{equation}
 c_d=\frac{\mu(d)\log d}{d}
 \label{eq:coefficient}
\end{equation}
on a squarefree divisor band.  Its reciprocal emitter is coupled across
divisors through common rational frequencies.  TPC-213 identified this common
physical source, and TPC-214 reduced the complete-period Gram to
reduced-denominator clusters \cite{WangTPC214}.  TPC-215 then proved that the
cluster Gram is at most an \(O((\log x)^2)\) multiple of the divisor
direct-sum energy \cite{WangTPC215}.

The next question is a narrower one: how large can that direct-sum row energy
be before finite-window cross frequencies are restored?  The answer here is a
source-locked envelope, not a cancellation theorem.  The fixed-prime integer
cutoff has a useful exact feature.  Since the shell is contained in \((Q,2Q]\)
and \(4Q<H\) for large \(x\), distinct admissible integers cannot occupy the
same residue modulo a fixed divisor.  This removes cross terms within one
prime row.  Cauchy is used only once, across the prime shell, and the resulting
divisor sum is elementary.

There is also a sharp warning.  The prime rows need not be orthogonal to one
another.  A finite exact fixture with primes congruent to one modulo five makes
all rows occupy the same two residues.  Thus the shell Cauchy step is a genuine
bottleneck of this structural proof.

\section{Literal source lock}

We retain the V46 scales exactly as frozen in the transition compiler
\cite{WangV46}:
\begin{equation}
 H=x^{21/32},\qquad Q=x^{1/3},\qquad
 Y_0=\frac{H}{4Q},\qquad U=x^{133/400}.
 \label{eq:scales}
\end{equation}
Let
\begin{equation}
 \Qx=\{q\text{ prime}:Q<q\le 2Q\},\qquad P=|\Qx|,
 \label{eq:shell}
\end{equation}
and
\begin{equation}
 \Dx=\{d\in\ZZ:Y_0<d\le U,\ \mu(d)^2=1\}.
 \label{eq:band}
\end{equation}
For sufficiently large \(x\), the source inequalities give \(U<Q\), so
every shell prime is invertible modulo every \(d\in\Dx\).  For a bounded
profile \(\psi\), set
\begin{equation}
 B_d(r)=\sum_{q\in\Qx}\sum_{0<|m|\le\lfloor dq/H\rfloor}
 \psi\!\left(\frac{Hm}{dq}\right)
 \one_{m\overline q\equiv r\pmod d}.
 \label{eq:emitter}
\end{equation}
Write \(M_\psi=\sup_t|\psi(t)|\) and let
\(\norm{v}_{2,d}^2=\sum_{r\bmod d}|v(r)|^2\).

The exact exponent ledger is
\begin{equation}
 \frac{H}{4Q}=\frac14x^{31/96}\longrightarrow\infty,
 \qquad \frac{U}{Q}=x^{-1/1200}\longrightarrow0.
 \label{eq:source-inequalities}
\end{equation}
Consequently, the proof below works in the asymptotic source range where
\(4Q<H\).  The factor four is important: the shell extends to \(2Q\), and
fixed-prime injectivity requires \(2q<H\), not merely \(2Q<H\).

\section{Fixed-prime injectivity}

For one prime define
\begin{equation}
 B_{d,q}(r)=\sum_{0<|m|\le n_{d,q}}
 \psi\!\left(\frac{Hm}{dq}\right)
 \one_{m\overline q\equiv r\pmod d},
 \qquad n_{d,q}=\left\lfloor\frac{dq}{H}\right\rfloor.
 \label{eq:fixed-row}
\end{equation}

\begin{theorem}[fixed-q no-collision]
\label{thm:no-collision}
For \(d\in\Dx\) and \(q\in\Qx\), distinct admissible integers \(m\) in
\eqref{eq:fixed-row} give distinct residues modulo \(d\).  Hence
\begin{equation}
 \norm{B_{d,q}}_{2,d}^2
 =\sum_{0<|m|\le n_{d,q}}
 \left|\psi\!\left(\frac{Hm}{dq}\right)\right|^2
 \le 2M_\psi^2\frac{dq}{H}.
 \label{eq:fixed-energy}
\end{equation}
\end{theorem}

\begin{proof}
Suppose \(m_1\ne m_2\) collide.  Since \(q\) is a unit modulo \(d\),
\(d\mid(m_1-m_2)\).  On the other hand, \(q\le2Q\) and \eqref{eq:source-inequalities}
give
\[
0<|m_1-m_2|\le2\left\lfloor\frac{dq}{H}\right\rfloor
 \le\frac{2dq}{H}<d.
\]
This is impossible.  Thus the atom supports are distinct, so the squared
norm is the sum of squared atom weights.  There are at most
\(2\lfloor dq/H\rfloor\le2dq/H\) atoms, proving \eqref{eq:fixed-energy}.
\end{proof}

\section{The direct-sum envelope}

\begin{theorem}[direct-sum row-energy envelope]
\label{thm:envelope}
Let
\begin{equation}
 \Edir=L\sum_{d\in\Dx}|c_d|^2\norm{B_d}_{2,d}^2,
 \label{eq:direct-energy}
\end{equation}
where \(L\) is a complete common period.  Then
\begin{equation}
 \frac{\Edir}{L}
 \le 16M_\psi^2\frac{Q^3}{H}
 \sum_{Y_0<d\le U}\frac{\mu(d)^2(\log d)^2}{d}
 \ll_{\psi}\frac{Q^3}{H}(\log U)^3.
 \label{eq:envelope}
\end{equation}
In particular,
\begin{equation}
 \frac{\Edir}{L}\ll_{\psi}x^{11/32}(\log x)^3,
 \label{eq:exponent}
\end{equation}
The right-hand scale is \(x^{11/32+o(1)}\), but no matching lower bound is
claimed.
\end{theorem}

\begin{proof}
Since \(B_d=\sum_{q\in\Qx}B_{d,q}\), Cauchy in the \(P\) shell coordinates
and Theorem~\ref{thm:no-collision} give
\begin{align}
 \norm{B_d}_{2,d}^2
 &\le P\sum_{q\in\Qx}\norm{B_{d,q}}_{2,d}^2 \\
 &\le \frac{2M_\psi^2Pd}{H}\sum_{q\in\Qx}q
 \le \frac{4M_\psi^2P^2dQ}{H}.
 \label{eq:shell-cauchy}
\end{align}
The interval \((Q,2Q]\) contains at most \(2Q\) integers for \(Q\ge1\),
so \(P\le2Q\).  Therefore
\[
 \norm{B_d}_{2,d}^2\le16M_\psi^2\frac{dQ^3}{H}.
\]
Using \(|c_d|^2=\mu(d)^2(\log d)^2/d^2\) in \eqref{eq:direct-energy}
proves the first inequality in \eqref{eq:envelope}.  Finally,
comparison with the integral of \((\log t)^2/t\) gives
\[
 \sum_{1\le d\le U}\frac{(\log d)^2}{d}\ll(\log(2U))^3.
\]
Since \(Q^3/H=x^{1-21/32}=x^{11/32}\) and \(\log(2U)=O(\log x)\),
the result follows.
\end{proof}

\begin{remark}[claim boundary]
The proof uses absolute values in both the fixed-row bound and shell Cauchy.
It does not use the signs of \(\mu(d)\), any prime cancellation, or the
four-packet signed reassembly.  Also, \(\Edir\) is a complete-period quantity.
The finite physical interval has an off-frequency Gram that is not controlled
by Theorem~\ref{thm:envelope}.
\end{remark}

\section{Exact aligned-shell adversary}

The shell Cauchy step cannot be replaced by a direct orthogonal sum.  Consider
the finite rational fixture
\begin{equation}
 d=5,\qquad H=500,\qquad Q_{\mathrm{scale}}=100,
 \qquad \{q\} = \{101,131,151,181\},
 \qquad \psi(t)=\frac{1}{(1+t^2)^2}.
 \label{eq:adversary}
\end{equation}
Each displayed \(q\) is prime and congruent to one modulo five.  Moreover,
\(\lfloor5q/500\rfloor=1\), so the only nonzero integers are \(m=\pm1\).
Since \(q^{-1}=1\pmod5\), every fixed-q row has support exactly
\(\{1,4\}\).  The exact rational certificate in the accompanying repository
records the individual rows, the combined row, and all norms without floating
point arithmetic.

For orientation, the exact values have decimal displays
\begin{center}
\begin{tabular}{@{}lr@{}}
\toprule
quantity & decimal display \\
\midrule
sum of individual row norms & 1.604830605122346 \\
combined row norm & 5.946998427183525 \\
combined/direct ratio & 3.705686075652919 \\
\bottomrule
\end{tabular}
\end{center}
The decimal values are only presentations of exact rationals.  The ratio is
strictly larger than one, so positive cross terms occur.  The fixture is a
finite structural adversary, not a model for an asymptotic prime shell.

\begin{proposition}[aligned-shell obstruction]
\label{prop:adversary}
No proof that uses only the individual fixed-q row norms may replace
\(\norm{\sum_qB_{d,q}}_2^2\) by \(\sum_q\norm{B_{d,q}}_2^2\) for the literal row
family.  In particular, free shell orthogonality is refuted in this scoped
structural sense.
\end{proposition}

\begin{proof}
The fixture \eqref{eq:adversary} has common support and an exact combined norm
strictly larger than the individual norm sum.  Therefore the proposed
orthogonal replacement fails on an admissible finite reciprocal-emitter
configuration.  The statement is scoped to the replacement rule; it does not
assert an asymptotic lower bound for the V46 object.
\end{proof}

\section{Route evaluation and open gates}

Theorem~\ref{thm:envelope} is the next quantitative envelope after TPC-215:
\begin{equation}
 \text{complete-period direct row energy}
 \le x^{11/32+o(1)}.
 \label{eq:route-advance}
\end{equation}
It removes the need to leave this direct-sum object completely unbounded, but
it does not provide a saving relative to a required target scale.  The exact
remaining gates are:
\begin{enumerate}[leftmargin=2em]
 \item attach the normalized complete-period envelope to the literal finite
       physical window and its off-frequency Gram;
 \item preserve the Mobius signs, prime-only shell, block tails, and four-packet
       reassembly in one theorem;
 \item pay the arithmetic \(L^2\), fixed-atom, and strict \(1/400\) ledgers.
\end{enumerate}
Route A is not applicable to this analytic twin-prime session.  The current
Route-B structural threshold passes, while arithmetic advance is
\texttt{NO}, \(L^2=\texttt{NONE}\), fixed-atom credit is zero, and full Gate B
is open.

\section{Conclusion}

The source inequality \(4Q<H\) turns each fixed-prime reciprocal emitter row
into an exact sum of distinct integer atoms.  Shell Cauchy and the elementary
Mobius-log divisor sum then give the complete-period envelope
\(L^{-1}\Edir\ll_\psi x^{11/32}(\log x)^3\).  The aligned rational fixture
shows why the shell Cauchy factor is a real bottleneck rather than a notation
artifact.  The next research step is a finite-window attachment that keeps this
alignment issue and the literal signed reassembly visible.

\bibliographystyle{plain}
\bibliography{references}

\end{document}
